\section{导读：为什么这期要当作一份季度讲义来读}

本期是“全球大模型季报”的跨年对谈。它不是普通人物访谈，而是一次面向 2026 年的 AI 产业、模型范式和投资主线复盘。张小珺负责把问题推进到商业、平台和创业层面，广密则把一组季度观察压缩成几个关键词：AI War、交替领先、基础模型像综合电商、OpenAI 与 Google 的双主线、Pre-training/RL/Online Learning 三种范式，以及“模型即产品，数据即模型”。

因此，本讲义不按原始字幕逐句展开，而按教学逻辑重组为六条主线。第一条是为什么“AI Bubble”不足以解释这一轮投入，AI 更像输不起的军备竞赛；第二条是 OpenAI 的收入账与 Google/OpenAI 的入口之争；第三条是 GPT、Claude、Gemini 的交替领先与基础模型公司的平台类比；第四条是第三范式 Online Learning；第五条是数据、Agentic Web、Neo Labs 与 Robotics；第六条是中美叙事、华人创业者和未来三到五年的长期变量。

\begin{importantbox}{本期核心命题}
广密的中心判断不是“没有泡沫”，而是“泡沫框架太窄”。AI 的投入同时被商业防守、国家战略、算力供给、人才竞争和下一范式预期驱动。它更像 AI War：没人确定短期账一定算得过来，但关键玩家都不敢下牌桌。
\end{importantbox}

\begin{knowledgebox}{视觉策略说明}
本视频没有 slides、白板、论文画面或产品演示。按照播客工作流，正文只使用封面做来源识别，不重复插入固定访谈画面；正文中的图均为自制概念图，用来承载结构关系，细节解释放在文字、表格和读图框里。
\end{knowledgebox}

\subsection{本章小结}

EP127 的价值在于把过去几期反复出现的模型、数据、Agent、算力和创业话题压缩成一个跨年判断：2026 年要看的不是单个 benchmark，而是范式是否变化、入口如何重分配、谁能把数据和产品做成闭环。

